\section{Discussion}

\label{sec:discussion}

\subsection{Implications for DeSci}

ZooCoord demonstrates that decentralized infrastructure can address concrete scientific workflow challenges, not merely provide alternative funding or credentialing mechanisms.
The key insight is that experiment management benefits from decentralization because trust requirements are fundamentally distributed: no single party should be trusted to verify its own work, yet centralized verification services introduce new trust dependencies.

\subsection{Limitations}

Several limitations qualify our findings.
First, the deployment population consists of researchers already engaged with the Zoo Network, introducing selection bias toward technology-comfortable participants.
Second, the 9-month deployment period is insufficient to assess long-term economic sustainability of the verification market.
Third, wet-lab experiments (e.g., conservation genetics) have inherently lower reproducibility than purely computational work, and our protocol can only verify the computational components.

\subsection{Scalability}

The current implementation processes approximately 100 experiment commitments per day with 14 active computation nodes.
Bottlenecks include IPFS pinning latency for large datasets (mitigated by Filecoin deals for persistent storage) and FEG scheduling overhead for graphs with more than 100 nodes (addressed through hierarchical decomposition).

% ============================================================
% 11. Conclusion
% ============================================================
